\section{Comparaison avec les ajustements des constantes fondamentales}
\label{sec:comparacion}

La comparaison externe évalue les conséquences d'expressions dont les états initiaux, les régions, les coefficients et les retenues ont été déterminés par la construction précédente. Pour chaque référence \(x_{\rm ref}\), d'incertitude-type \(u(x_{\rm ref})\), on indique le résidu normalisé
\[
 z=\frac{x_{\rm calc}-x_{\rm ref}}{u(x_{\rm ref})}.
\]
Ce rapport exprime la différence par rapport à chaque centre tabulé en unités de son incertitude-type. Une interprétation probabiliste exige de spécifier le modèle statistique ; la comparaison conjointe des ajustements CODATA exige en outre leurs corrélations et les conditions d'indépendance utilisées.

Les expressions utilisées dans cette comparaison donnent
\begin{align}
 \alpha&=0.007297352569283800997285105472380662999\ldots,\notag\\
 \frac{m_ec^2}{\mathrm{MeV}}&=0.510998950690000808608\ldots.
 \label{eq:valores-comparacion}
\end{align}
La seconde valeur correspond à l'expression électronique corrigée, avec le facteur complet de \(\Delta_4\) dans la normalisation basale et dans l'exposant de retour. La composition du terme basal avec la mémoire de retour produit cette valeur avant l'introduction de la masse tabulée comme référence de comparaison.

\subsection{Exactitude, incertitude et réalisation dimensionnelle}

La comparaison articule trois grandeurs aux fonctions distinctes. L'exactitude algébrique correspond aux identités démontrées dans leurs domaines. L'erreur d'évaluation numérique est contrôlée au moyen de la précision du calcul. L'incertitude expérimentale quantifie la connaissance de la grandeur mesurée selon l'ajustement utilisé et conserve ce sens lorsqu'on étend l'évaluation décimale. Les chiffres supplémentaires de \eqref{eq:valores-comparacion} décrivent l'expression mathématique ; l'information expérimentale est représentée par les centres, les incertitudes et les corrélations de l'ajustement.

La réalisation dimensionnelle de la section~\ref{sec:electron} spécifie l'application du caractère électronique à l'énergie et à la masse, avec son échelle et sa normalisation. Cette application détermine la grandeur et les unités sur lesquelles porte la comparaison. La construction structurelle établit la valeur évaluée et ses relations ; la comparaison métrologique examine sa correspondance avec la grandeur physique ainsi identifiée. L'ordre causal conserve APP, TRIT et TPK comme antécédents des expressions évaluées, suivies de la réalisation dimensionnelle et des valeurs de référence.

\subsection{CODATA 2018}

L'ajustement de 2018 contient les valeurs centrales et les incertitudes suivantes\cite{codata2018} :
\begin{center}
\begin{tabular}{@{}lcccr@{}}\toprule
Grandeur&Valeur HMT&Référence&\shortstack{Incertitude\\type}&\(z\)\\\midrule
\(\alpha\)&\(0.0072973525692838\ldots\)&\(0.0072973525693\)&\(1.1\times10^{-12}\)&\(-0.01473\)\\
\(m_ec^2/\mathrm{MeV}\)&\(0.5109989506900008\ldots\)&\(0.51099895000\)&\(1.5\times10^{-10}\)&\(+4.60001\)\\\bottomrule
\end{tabular}
\end{center}
L'expression de \(\alpha\) coïncide avec la valeur centrale de 2018 après arrondi aux chiffres affichés. La différence, d'environ \(-1.62\times10^{-14}\), représente environ \(0.0147\) fois l'incertitude-type imprimée. L'expression électronique corrigée dépasse le centre de 2018 d'environ \(6.90\times10^{-10}\,\mathrm{MeV}\), différence qui persiste à l'arrondi indiqué. Chaque ligne conserve ainsi son accord décimal ou son écart quantifié.

\subsubsection{Comparaison de la fenêtre dodécaphasique avec l'inverse du centre}
\label{sec:contraste-ventana-reciproco}

La fenêtre finie \(\alpha_{12}^{(0)}\) de
\eqref{eq:alpha-ventana-cero}, obtenue avec \(c_{13}=0\), admet en outre
la comparaison historique avec l'inverse du centre imprimé de la
constante inverse. Définissons
\[
 R_{2018}:=137.035999084,\qquad
 \alpha_{\mathrm{rec},2018}:=R_{2018}^{-1}.
\]
Dans cette opération, on évalue l'inverse du rationnel décimal
\(R_{2018}\). Cette référence se distingue du centre tabulé
\(0.0072973525693\) utilisé dans le tableau précédent.

\begin{center}
\begin{tabular}{@{}ll@{}}
\toprule
Expression & Évaluation arithmétique\\
\midrule
\(\alpha_{12}^{(0)}\) &
\(0.007297352569283800997285105472380663\)\\
\(\alpha_{\mathrm{rec},2018}\) &
\(0.007297352569283800997285105472380663567972560701\ldots\)\\
\(\alpha_{12}^{(0)}-\alpha_{\mathrm{rec},2018}\) &
\(\simeq-5.679725607012965\times10^{-37}\)\\
\((\alpha_{12}^{(0)})^{-1}-R_{2018}\) &
\(\simeq+1.066588006664435\times10^{-32}\)\\
\(\bigl((\alpha_{12}^{(0)})^{-1}-R_{2018}\bigr)/R_{2018}\) &
\(\simeq+7.783268730800000\times10^{-35}\)\\
\bottomrule
\end{tabular}
\end{center}

La relation décimale précise est
\[
 \alpha_{12}^{(0)}
 =10^{-36}\left\lfloor10^{36}\alpha_{\mathrm{rec},2018}\right\rfloor.
\]
Par conséquent, la coïncidence historique correspond à la troncature à
trente-six positions décimales de l'inverse ; son arrondi à cette
position se termine par \(664\), tandis que la fenêtre se termine par \(663\).
Les différences du tableau quantifient la comparaison entre les deux
nombres. La section infinie \(\alpha_A\) conserve sa propre normalisation
compatible et sa propre frontière, distinctes de celles de cette fenêtre finie.

L'ajustement exprime la constante inverse sous la forme
\(137.035999084(21)\), avec une incertitude-type
\(u(R_{2018})=2.1\times10^{-8}\)\cite{codata2018}.
Les décimales obtenues en inversant son centre imprimé appartiennent à
l'évaluation arithmétique de ce rationnel ; l'incertitude expérimentale
reste celle indiquée par l'ajustement. La comparaison précédente avec le
centre tabulé de \(\alpha\) maintient sa référence et son résidu
normalisé propres.


\subsection{CODATA 2022}

L'ajustement 2022 modifie les valeurs centrales\cite{codata2022} :
\begin{center}
\begin{tabular}{@{}lcccr@{}}\toprule
Grandeur&Valeur HMT&Référence&\shortstack{Incertitude\\type}&\(z\)\\\midrule
\(\alpha\)&\(0.0072973525692838\ldots\)&\(0.0072973525643\)&\(1.1\times10^{-12}\)&\(+4.53073\)\\
\(m_ec^2/\mathrm{MeV}\)&\(0.5109989506900008\ldots\)&\(0.51099895069\)&\(1.6\times10^{-10}\)&\(+0.00000505\)\\\bottomrule
\end{tabular}
\end{center}
Dans cette édition, l'expression électronique coïncide avec la valeur centrale tabulée à l'arrondi près, tandis que \(\alpha\) s'en écarte d'environ \(4.53\) incertitudes-types imprimées. La comparaison de l'exactitude physique des ajustements repose sur leurs données expérimentales, leurs corrélations et leurs hypothèses. La date et la concordance décimale avec une expression identifient des propriétés de la comparaison ; l'évaluation de la proximité avec la valeur physique exige cette analyse expérimentale complète.

\Needspace{29\baselineskip}
\subsection{Confrontation de la constante de Planck réduite}
\label{sec:revision-planck-contraste}

Le Système international en vigueur depuis le 20 mai 2019 fixe
exactement \(h=6.62607015\times10^{-34}\,\mathrm{J\,s}\)
\cite{bipm2018}. Sa constante réduite est donc
\[
 \hbar_{\mathrm{SI}}=\frac{6.62607015}{2\pi}\,10^{-34}\,
 \mathrm{J\,s}
 =1.0545718176461563912\ldots\times10^{-34}\,\mathrm{J\,s}.
\]
Les points de suspension correspondent au développement d'une expression
exacte, et non à une incertitude expérimentale de \(h\) dans le SI actuel.
Cette valeur de référence est utilisée ici après l'évaluation des
lecteurs de HMT ; sa fonction est comparative.

Les deux coordonnées du retour de l'article sont
\begin{align*}
 H_5&=1.0545718176461544696\ldots,\\
 \eta_{\mathrm{ret}}&=1.0545718169150390611\ldots.
\end{align*}
Dans l'identification de \(\mathcal U_S\) à \(\mathrm{J\,s}\), la
différence relative de la section antérieure par rapport à
\(\hbar_{\mathrm{SI}}\) est d'environ \(-1.82\times10^{-15}\).
Celle de la section postérieure est d'environ \(-6.93\times10^{-10}\).
La seconde différence contient l'effet du retour utilisé par le
correcteur électronique ; elle n'équivaut pas à l'erreur d'évaluation de la
première section.

La précision de ces comparaisons permet de reconnaître la proximité
quantitative de la construction de \(H_5\) et, simultanément, de distinguer
proximité et identité exacte. Dans le présent manuscrit, les expressions
présentées ne fournissent pas une égalité exacte avec \(h/(2\pi)\).
Le résultat démontré sur \(\Sigma_H\) est sa décade \(-34\) ;
le résultat analytique sur \(H_5\) est l'évaluation du caractère
déclaré. L'identification physique est examinée dans cette coordinatisation et avec les
différences explicites que nous venons de donner. La valeur SI n'intervient
ni dans la définition de la norme déterminantale ni dans la récurrence régionale.

